\Needspace{15\baselineskip}
\section{子流形与局部流习题及详细解答}
习题5.3、5.4以梁书上册第二版2.2.2的文内命题和选读说明为出发点，增加坐标计算与定义域核查；并非原书章末题号。其余六题为原创补充。来源位置与双页码见\texttt{references/EXERCISES.md}，全部解答由本工程编写。

\Needspace{10\baselineskip}
\subsection*{习题5.1：双约束与切空间}
设$F:\R^3\to\R^2$，$F(x,y,z)=(x^2+y^2+z^2,z)$。判断$(1,c)$何时为非空纤维的正则值，并求纤维及其切空间。解释$c=\pm1$和$|c|>1$。
\begin{proof}[解答]
纤维满足$z=c$、$x^2+y^2=1-c^2$。当$|c|<1$时为水平圆，$DF$的两行是$(2x,2y,2c)$与$(0,0,1)$，因$(x,y)\ne(0,0)$而线性无关，故为正则值。切空间为
\[
 \{(a,b,d):d=0,\ xa+yb=0\}
 =\operatorname{span}\{(-y,x,0)\}.
\]
当$c=\pm1$时纤维是单点$(0,0,c)$，其本身是零维子流形，但$DF$秩为一，值不是正则值，不能用$m-n=1$推断其维数。$|c|>1$时纤维为空，按定义仍是正则值。可见“正则值”不保证纤维非空，奇异值也不强迫纤维不是子流形。
\end{proof}

\Needspace{9\baselineskip}
\subsection*{习题5.2：自交浸入与紧源}
证明$\gamma:S^1\to\R^2$、$\gamma([t])=(\sin t,\sin2t)$为浸入而非嵌入。说明为什么这不违背紧源命题；再证明其限制到充分小的参数弧上为嵌入。
\begin{proof}[解答]
导数为$(\cos t,2\cos2t)$；若$\cos t=0$，则$\cos2t=-1$，所以两分量不会同时为零，故为浸入。$[0]\ne[\pi]$却同映到$(0,0)$，所以不单射，也就不是嵌入。命题\ref{prop:compact-embedding}同时要求单射，不能只用紧致性。对任意$t_0$，至少一个分量的导数非零，该分量在小参数弧上由逆映射定理给出局部参数，因而$\gamma$在该弧上单射且到像的逆连续、光滑。两个分支在环境中相遇，不妨碍源上每个分支独立地局部嵌入。
\end{proof}

\Needspace{10\baselineskip}
\subsection*{习题5.3：对易子、坐标基与函数系数}
依据梁书印刷页28定义10及定理2-2-7，证明坐标场$\partial_x,\partial_y$的括号为零。在$\R^2$取$X=\partial_x,Y=x\partial_y$，计算$[X,Y]$，并在原点比较$Y$与零场。向量场的李括号能否作为$T_0\R^2$上的双线性运算？
\begin{proof}[解答]
对局部光滑$f$，$[\partial_x,\partial_y]f=\partial_x\partial_yf-\partial_y\partial_xf=0$。一般式则给$[\partial_x,x\partial_y]=\partial_x(x)\partial_y=\partial_y$。原点处$Y_0=0$，但$[X,Y]_0=\partial_y|_0$，而$[X,0]_0=0$。所以括号的值还读取向量场邻域内的变化，不能只由两个点向量定义；若强行定义为点态双线性运算，与零向量的括号应为零，立即矛盾。这里的双线性是对实常数的双线性，不是对光滑函数的双线性。
\end{proof}

\Needspace{10\baselineskip}
\subsection*{习题5.4：最大积分曲线的时间区间}
参照梁书印刷页28--29定理2-2-8及选读2-2-3的局部唯一性与不可延拓说明，求$X=x^2\partial_x$各初值的极大积分曲线。限制到$M=(-1,1)$时，从$p=0$与$p=1/2$出发的极大时间区间分别是什么？
\begin{proof}[解答]
$p=0$时由唯一性得到恒零解，区间为$\R$。$p\ne0$时在解不为零的区间积分$(1/x)'=-1$，得$x(t)=p/(1-pt)$。正$p$的含零区间为$(-\infty,1/p)$，负$p$为$(1/p,\infty)$；有限端点处$|x|\to\infty$，不能连续延拓为实线值解。到无穷远的另一端没有有限时间端点，故这些区间极大。
在$M=(-1,1)$上，$p=0$仍为恒零全局解。$p=1/2$时$x(t)=1/(2-t)$，要求$-1<x(t)<1$并留在含零的连续分支，得到$t<1$，故区间为$(-\infty,1)$。当$t\uparrow1$时趋向被删去的边界点$1$，虽未在$\R$中爆破，却不能继续作为$M$值解。这说明极大区间同时依赖方程和所在流形。
\end{proof}

\Needspace{10\baselineskip}
\subsection*{习题5.5：显式流盒}
对$X=\partial_x+x\partial_y$求流，构造全平面上的流盒坐标，并证明这些坐标确实使$X$变为一个坐标场。是否由此可知该换坐标保持长度？
\begin{proof}[解答]
从$(x_0,y_0)$出发，$x=x_0+t$，积分$\dot y=x_0+t$得
\[
 \Phi_t(x_0,y_0)=(x_0+t,y_0+x_0t+t^2/2).
\]
令$s=x,z=y-x^2/2$，其逆为$(x,y)=(s,z+s^2/2)$，两者全局光滑，Jacobi行列式恒为一。$X(s)=1$、$X(z)=x-x=0$，所以新坐标中$X=\partial_s$。然而在$(s,z)$的坐标下，欧氏度量为$\dd s^2+(\dd z+s\dd s)^2$，并非$\dd s^2+\dd z^2$，所以换坐标不保持以标准分量计算的长度。流盒结论属于光滑结构，未给出等距性。
\end{proof}

\Needspace{10\baselineskip}
\subsection*{习题5.6：可积平面场与第一积分}
对$D=\operatorname{span}(\partial_x+y\partial_z,\partial_y+x\partial_z)$，求非恒定第一积分和经过$(a,b,c)$的局部积分片；验证两个生成场的流交换。
\begin{proof}[解答]
置$w=z-xy$，则$\dd w=(-y,-x,1)$处处非零，且两生成场均使$w$的导数为零。两场逐点独立、核为二维，故$D=\ker\dd w$。经过指定点的片为$z-xy=c-ab$，切空间等于该核。
第一生成场的流为$\Phi_s(x,y,z)=(x+s,y,z+sy)$，第二场的流为$\Psi_t(x,y,z)=(x,y+t,z+tx)$。两个复合均为
\[
 (x+s,y+t,z+sy+tx+st),
\]
故全局交换。第一积分把局部曲面方程与流的相容性联系起来，但不因此将所有分布都视为某个全局函数的核。
\end{proof}

\Needspace{10\baselineskip}
\subsection*{习题5.7：不可积的混合偏导障碍}
对$E=\operatorname{span}(\partial_x,\partial_y+x\partial_z)$，除李括号判别外，直接证明不存在二维积分曲面。说明为何局部曲面可以先假设为图像$z=h(x,y)$。
\begin{proof}[解答]
若有积分浸入，其切平面投影到$(x,y)$平面是线性同构：两生成向量的投影是标准基。因此对积分浸入复合该投影应用逆映射定理，局部就能把曲面写成图像。图像切空间为$\operatorname{span}(\partial_x+h_x\partial_z,\partial_y+h_y\partial_z)$，与$E$相同且$x,y$投影固定时，唯一提升给出$h_x=0$、$h_y=x$。于是$h_{xy}=0$而$h_{yx}=1$，违反光滑函数混合偏导交换，矛盾。这个证明同时解释了为什么不能凭图像直觉省略局部投影可逆这一步。
\end{proof}

\Needspace{10\baselineskip}
\subsection*{习题5.8：对合标架未必是坐标标架}
在$\R^2$取$Z_1=\partial_x,Z_2=e^x\partial_y$。证明其张成的分布对合，但不存在局部坐标使这两个指定向量场同时成为两个坐标基向量。给出一个光滑标架变换，将它们改成坐标标架。
\begin{proof}[解答]
两场处处独立，张成$T\R^2$，该分布当然对合；也可计算$[Z_1,Z_2]=e^x\partial_y=Z_2$并用标架判别。若两指定场同时为坐标基，则它们的括号应为零，与该计算矛盾。改取$Y_1=Z_1,Y_2=e^{-x}Z_2=\partial_y$，得到原坐标的基，二者交换。Frobenius只保证分布有一个坐标标架，并不承诺输入的任意标架都可直接成为坐标基；正文中的归一化步骤正是为解决这个差别。
\end{proof}
